\section*{الفصل \ref{ch:b1:fluid-envelopes-structure} --- بنية الغلاف الجوي والمحيط}
% القيم: calc/ch26.py

\begin{solution}{exo:b1:fluid-envelopes-structure:1}
(أ)~$H = R^*T/(M_0 g_0) = 8314.32 \times 250/(28.9644 \times 9.80665) =
\qty{7.32}{km}$. (ب)~$e^{10/7.32} = 3.92$: ينخفض الضغط إلى نحو الربع. (ج)~$H$ يتناسب مع الحرارة: عند \qty{216.65}{K} يكون \qty{6.34}{km}، مقابل \qty{8.43}{km} عند \qty{288.15}{K}. فالهواء البارد أكثف عند ضغط معيّن، فينخفض الضغط أسرع مع الارتفاع.
\end{solution}

\begin{solution}{exo:b1:fluid-envelopes-structure:2}
مع $\Gamma_d = \qty{9.76}{K/km}$: \qtyrange{0}{0.3}{km}، $11.0 > 9.76$،
غير مستقرّة؛ \qtyrange{0.3}{1.2}{km}، $9.7$، محايدة في حدود دقّة البيانات؛ \qtyrange{1.2}{1.5}{km}، انعكاس؛ \qtyrange{1.5}{8}{km}، مستقرّة؛
فوق \qty{8}{km}، متساوية الحرارة ومستقرّة. الانعكاس هو الأشدّ استقرارًا، لأن $\Gamma_d - \Gamma = 13.8$ أكبر ما يكون هناك.
\end{solution}

\begin{solution}{exo:b1:fluid-envelopes-structure:3}
(أ)~$\sigma = 1026.95 - 1000 = \qty{26.95}{kg.m^{-3}}$. (ب)~B، الأخفّ، يعلو A. (ج)~جميع كثافات ماء البحر تبدأ بالرقمين نفسيهما، 10؛ و$\sigma$ يحتفظ فقط بالأرقام التي تتغيّر من ماء إلى آخر.
\end{solution}

\begin{solution}{exo:b1:fluid-envelopes-structure:4}
(أ)~$S = 0.25 \times 35.40 + 0.75 \times 34.30 = 34.575$؛ $\theta = 0.25 \times
15.0 + 0.75 \times 4.0 = \qty{6.75}{\celsius}$. (ب)~على القطعة PQ، على ربع المسافة من Q نحو P (\cref{prop:b1:fluid-envelopes-structure:mixing}).
(ج)~خلائط P وQ وR تملأ المثلث PQR. بكتابة $(34.60, 8.0) = a\,\mathrm P + b\,\mathrm Q + c\,\mathrm R$ مع $a + b + c = 1$
نجد $a = 0.345$، و$b = 0.788$، و$c = -0.133$: نسبة سالبة من R، فالنقطة تقع خارج المثلث، ولا يمكن لخليط من الثلاثة أن ينتجها.
\end{solution}

\begin{solution}{exo:b1:fluid-envelopes-structure:5}
(أ)~$\theta = 253.15 \times 2^{2/7} = \qty{308.6}{K}$. (ب)~عند \qty{1000}{hPa}،
$\theta = T = \qty{298.15}{K}$. تزداد \omterm{def:b1:fluid-envelopes-structure:potential-temperature}{درجة الحرارة الكامنة} نحو الأعلى، من 298.2 إلى \qty{308.6}{K}: الطبقة مستقرّة للهواء الجاف. (ج)~كان سيصل بدرجته الكامنة، \qty{308.6}{K}، أي \qty{35.4}{\celsius}، أدفأ بعشر درجات من الهواء الموجود هناك.
\end{solution}

\begin{solution}{exo:b1:fluid-envelopes-structure:6}
(أ)~$T = 288.15 - 6.5 \times 5 = \qty{255.65}{K}$ و$p = \num{101325} \times
(255.65/288.15)^{5.256} = \qty{54020}{Pa}$، أي \qty{540.2}{hPa}. (ب)~$H =
\qty{8.43}{km}$ و$p = \num{101325}\,e^{-5/8.43} = \qty{56010}{Pa}$. الضغط الحقيقي (المعياري) أدنى: فالهواء هناك أبرد من \qty{288.15}{K}، ومن ثمّ أكثف، وينخفض الضغط أسرع مع الارتفاع. (ج)~$p = \num{22632}\,
e^{-9/6.34} = \qty{5475}{Pa}$، القيمة الواردة في
\cref{tab:b1:fluid-envelopes-structure:layers}.
\end{solution}

\begin{solution}{exo:b1:fluid-envelopes-structure:7}
(أ)~$\beta\,\Delta S = \alpha\,\Delta T$ يعطي $\Delta S = 1.67/7.61 = 0.22$.
(ب)~$\Delta\rho = 1026.95 \times (-1.67\times10^{-4} \times 2 + 7.61\times10^{-4}
\times 0.3) = \qty{-0.11}{kg.m^{-3}}$: يصير الماء أخفّ.
(ج)~$\Delta\rho = 1028.11 \times (-0.53\times10^{-4} \times 2 + 7.85\times10^{-4}
\times 0.3) = \qty{+0.13}{kg.m^{-3}}$: التغيّرات نفسها تجعله أثقل الآن، لأن الاحترار بالكاد يغيّر كثافة الماء قرب درجة التجمّد.
\end{solution}

\begin{solution}{exo:b1:fluid-envelopes-structure:8}
(أ)~$N^2 = (9.81/260) \times (9.76 - 4)\times10^{-3} = \qty{2.17e-4}{s^{-2}}$،
$N = \qty{0.0147}{s^{-1}}$، والزمن $2\pi/N = \qty{426}{s}$، أي نحو سبع
دقائق. (ب)~$N^2 = (9.81/270) \times 14.76\times10^{-3} = \qty{5.36e-4}{s^{-2}}$،
$N = \qty{0.0232}{s^{-1}}$، والزمن \qty{271}{s}. (ج)~$N$ يتّجه إلى الصفر، والزمن يزداد بلا حدّ: قوة الاسترجاع تنعدم في طبقة محايدة.
\end{solution}

\begin{solution}{exo:b1:fluid-envelopes-structure:9}
(أ)~هو أثقل، $\sigma_\theta = 27.20$ مقابل 26.30: فAAIW أعذب لكنه أبرد بكثير. (ب)~$\Delta\sigma = 1026.95 \times [-1.67\times10^{-4} \times
(4.31 - 15.43) + 7.61\times10^{-4} \times (34.30 - 35.54)] = 1.907 - 0.969 =
0.94$، مقابل $0.90$ بالضبط. (ج)~الحرارة: البرودة تضيف
\qty{1.91}{kg.m^{-3}}، والعذوبة تنزع \qty{0.97}{kg.m^{-3}}.
\end{solution}

\begin{solution}{exo:b1:fluid-envelopes-structure:10}
(أ)~بموجب \cref{thm:b1:fluid-envelopes-structure:stability}، $(g/\theta)\,\mathrm
d\theta/\mathrm dz = (g/T)(\Gamma_d - \Gamma)$؛ و$\mathrm d\theta/\mathrm dz = 0$
يعطي $\Gamma = \Gamma_d$. (ب)~$T = T_s - \Gamma_d z$، طبقة ذات تدرّج ثابت
$L = \Gamma_d$: بموجب \cref{prop:b1:fluid-envelopes-structure:layer-pressure}، $p =
p_s (T/T_s)^{gM/(R\Gamma_d)}$، و$gM/(R\Gamma_d) = gMc_p/(Rg) = c_pM/R$
($= 3.5$ للهواء الجاف). (ج)~ينعدم $p$ حيث $T = 0$، عند $z = T_s/\Gamma_d =
300/9.76 = \qty{30.7}{km}$. الطبقات الحقيقية ذات \omterm{def:b1:fluid-envelopes-structure:potential-temperature}{درجة الحرارة الكامنة} المنتظمة سمكها بضعة كيلومترات على الأكثر: إذ يحدّها هواء مستقرّ، كما يحدّ الستراتوسفيرُ التروبوسفيرَ.
\end{solution}

\begin{solution}{exo:b1:fluid-envelopes-structure:11}
(أ)~$S = (33.787 + 36.050)/2 = 34.919$، $\theta = \qty{8.0}{\celsius}$.
(ب)~مع معادلة حالة خطّية، يكون $\sigma$ دالّة تآلفية في $S$ و$\theta$، فيكون $\sigma$ للخليط متوسط الاثنين، أي 27.00.
(ج)~$\alpha$ يساوي \qty{0.75e-4}{K^{-1}} عند \qty{2}{\celsius} و
\qty{2.07e-4}{K^{-1}} عند \qty{14}{\celsius}، بينما تتغيّر $\beta$ قليلًا
(\num{7.80e-4} و\num{7.53e-4}). تنخفض الكثافة أسرع فأسرع مع ارتفاع الحرارة: إنها دالّة مقعّرة للحرارة. وعلى خط الخلط، يكون الخليط عند الحرارة المتوسطة، حيث تقع الكثافة فوق الوتر الواصل بين أصليه، أي فوق 27.00. (د)~الزيادة تأتي من الانحناء، وهو حدّ من الرتبة الثانية: ينمو بمربّع فرق الحرارة بين أصليه.
\end{solution}

\begin{solution}{exo:b1:fluid-envelopes-structure:12}
(أ)~فرق الضغط هو $5604 - 5500 = \qty{104}{dbar}$: يبرد الطرد بمقدار
$0.129 \times 0.104 = \qty{0.013}{\celsius}$، إلى \qty{0.718}{\celsius}.
(ب)~الماء المحيط به عند \qty{5400}{m} عند \qty{0.730}{\celsius}، أدفأ: فالطرد، بملوحة تكاد تكون نفسها، أثقل من محيطه فيهبط ويعود. الطبقة مستقرّة، مع أن الحرارة المقاسة تزداد نحو الأسفل؛ إذ تتناقص \omterm{def:b1:fluid-envelopes-structure:potential-temperature}{درجة الحرارة الكامنة} نحو الأسفل. (ج)~$\theta \approx
0.7313 - 0.086 \times 5.604 = \qty{0.249}{\celsius}$، في حدود جزء من ألف من الدرجة من القيمة الدقيقة.
\end{solution}

\begin{solution}{pb:b1:fluid-envelopes-structure:1}
\begin{enumerate}
  \item الحرارة شبه منتظمة حتى \qty{50}{m}، وتهبط بسرعة حتى
        \qty{300}{m}، ثم أبطأ حتى \qty{1000}{m}، ثم شبه منتظمة تحتها.
        وتتناقص الملوحة من 35.80 عند السطح إلى حدّ أدنى قدره 34.32 عند
        \qty{800}{m}، ثم تزداد حتى 34.90 عند \qty{3000}{m}.
  \item تُبلغ $22.40 - 0.5 = \qty{21.90}{\celsius}$ بين \qty{50}{m}
        (22.10) و\qty{75}{m} (19.20): $50 + 25 \times 0.20/2.90 = \qty{52}{m}$.
  \item $(22.10 - 13.00)/250 = \qty{0.0364}{\celsius/m}$، أي \qty{3.6}{\celsius}
        لكل \qty{100}{m}. يمتد \omterm{def:b1:fluid-envelopes-structure:ocean-layers}{الترموكلاين} من قاعدة \omterm{def:b1:fluid-envelopes-structure:ocean-layers}{الطبقة
        المختلطة}، قرب \qty{50}{m}، حتى نحو \qty{800}{m}؛ وهو أحدّ ما يكون بين
        \qty{50}{m} و\qty{75}{m} (\qty{11.6}{\celsius} لكل \qty{100}{m}).
  \item $\sigma = 1026.95 \times [1 - 1.67\times10^{-4} \times 12.4 +
        7.61\times10^{-4} \times 0.80] - 1000 = 25.45$، وهو أعلى ب\qty{0.74}{kg.m^{-3}} من القيمة الدقيقة.
        ماء السطح بعيد عن الحالة المرجعية، و$\alpha$
        أكبر بكثير عند \qty{22}{\celsius} منه عند \qty{10}{\celsius}: فالشكل
        الخطّي يقلّل من تأثير الخفّة الناتجة عن الدفء.
  \item الحدّ الأدنى للملوحة عند \qty{800}{m}، 34.32 عند \qty{4.8}{\celsius}، هو ماء الأطلسي المتوسّط القطبي الجنوبي، المتكوّن تحت الرياح الغربية المطيرة في المحيط الجنوبي. والحدّ الأقصى عند \qty{3000}{m}، 34.90 عند
        \qty{2.5}{\celsius}، هو ماء الأطلسي الشمالي العميق، المتكوّن في بحر غرينلاند ولابرادور.
  \item $(2.50 - 2.26)/3.042 = \qty{0.079}{\celsius}$ لكل \qty{1000}{dbar}.
  \item $\sigma_\theta$ يزداد مع العمق: فالعمود مستقرّ. \omterm{def:b1:fluid-envelopes-structure:ocean-layers}{والبيكنوكلاين} هو حيث تزداد $\sigma_\theta$ أسرع ما يكون، في بضع مئات الأمتار العليا، مع \omterm{def:b1:fluid-envelopes-structure:ocean-layers}{الترموكلاين}: \num{0.82} بين \qty{200}{m} و\qty{800}{m}، و\num{0.70} فقط عبر \qty{2200}{m} تحتها.
  \item $S = (35.45 + 34.32)/2 = 34.885$؛ $\theta = (14.90 + 4.80)/2 =
        \qty{9.85}{\celsius}$.
  \item C: 26.461؛ I: 27.310؛ الخليط: 26.886، وهو بالضبط متوسّط الاثنين (26.886)، لأن الشكل الخطّي تآلفي.
  \item متوسط الأصلين هو $(26.342 + 27.159)/2 = 26.751$؛ والخليط، عند 26.888، أثقل ب\qty{0.14}{kg.m^{-3}}. وهذا هو التكثّف الخلطي.
  \item بالشكل الخطّي، يعطي القياس 26.61 عند \qty{300}{m}، و\num{26.88} عند
        \qty{500}{m}، و\num{27.31} عند \qty{800}{m}: فكثافة الخليط تقع قرب
        \qty{500}{m}. والماء هناك، $(34.82, \qty{9.6}{\celsius})$، قريب جدًّا من الخليط $(34.885, \qty{9.85}{\celsius})$: فالقياس بين الماء المركزي وAAIW قريب بحدّ ذاته من خطّ خلط بينهما.
  \item $S = 0.25 \times 35.45 + 0.75 \times 34.32 = 34.60$، $\theta = 0.25
        \times 14.90 + 0.75 \times 4.80 = \qty{7.33}{\celsius}$. وباستيفاء القياس بين \qty{500}{m} و\qty{800}{m} نجد
        $(34.57, \qty{7.2}{\celsius})$ قرب \qty{650}{m}: قريبة جدًّا منه.
  \item بوحدة \si{K/km}: A--B 10.0؛ B--C 9.7؛ C--D $-5.5$؛ D--E 6.29؛ E--F 6.50؛
        F--G 6.57؛ G--H 0.33؛ H--I $-0.50$؛ I--J $-0.83$.
  \item A--B غير مستقرّة؛ B--C محايدة؛ C--D انعكاس؛ D--E وE--F وF--G مستقرّة؛
        G--H شبه متساوية الحرارة ومستقرّة؛ H--I وI--J انعكاسان. والتروبوبوز عند G، \qty{12.5}{km} و\qty{-59.0}{\celsius}، حيث تكفّ الحرارة عن الانخفاض.
  \item متوسط الحرارة هو \qty{214.65}{K}: $H = 287 \times 214.65/9.81 =
        \qty{6.28}{km}$، و$147\,e^{-4/6.28} = \qty{77.7}{hPa}$، وهو متّسق مع القيمة المقاسة، 78.
  \item الهواء عند A، \qty{22.0}{\celsius}، أبرد من البحر: فتتدفّق الحرارة من البحر إلى الهواء، الذي يُسخَّن من الأسفل. والطبقة A--B، بمعدّل انخفاض أكبر من $\Gamma_d$، هي النتيجة.
  \item عند B: $22.0 - 9.76 \times 0.5 = \qty{17.1}{\celsius}$، مقابل 17.0؛ وعند
        C: \qty{7.4}{\celsius}، مقابل 7.3؛ وعند D: \qty{3.5}{\celsius}، مقابل
        9.5. فالطرد أدفأ قليلًا من الهواء حتى C، وأبرد كثيرًا عند D.
  \item عند قاع الانعكاس، قرب \qty{1.5}{km}: إذ يدخل الانعكاس، يصير الطرد أبرد، ومن ثمّ أكثف، من الهواء المحيط به، فيتوقّف صعوده.
  \item تصعد الطرود المنطلقة من الطبقة السطحية غير المستقرّة عبرها وتخلطها حتى تصير \omterm{def:b1:fluid-envelopes-structure:potential-temperature}{درجة الحرارة الكامنة} منتظمة، فيصبح معدّل انخفاضها $\Gamma_d$؛ والانعكاس فوقها يوقف الخلط.
  \item بالكلفن: A 294.15؛ B 293.99؛ C 294.08؛ D 300.51؛ E 312.65؛ F 329.29؛ G
        346.28؛ H 369.50؛ I 446.99؛ J 600.91.
  \item $\theta$ يتناقص من A إلى B (غير مستقرّة)، ويكاد يكون ثابتًا من B إلى C (محايدة)، ويتزايد فوق ذلك (مستقرّة)، وأسرع ما يكون من C إلى D وفوق G: التصنيف نفسه كما في السؤال~14.
  \item E--F: $(329.29 - 312.65)/4 = \qty{4.2}{K/km}$؛ وH--I: $(446.99 -
        369.50)/4 = \qty{19.4}{K/km}$، أي نحو خمسة أضعاف.
  \item مع المتوسط $\theta = \qty{408.2}{K}$: $N^2 = (9.81/408.2) \times
        0.01937 = \qty{4.66e-4}{s^{-2}}$، و$N = \qty{0.0216}{s^{-1}}$، والزمن
        \qty{291}{s}، أي نحو خمس دقائق.
  \item $N^2 = (9.81/1027) \times (26.35 - 24.79)/150 = \qty{9.9e-5}{s^{-2}}$،
        $N = \qty{0.0100}{s^{-1}}$، والزمن \qty{630}{s}. فالستراتوسفير السفلى أشدّ استقرارًا، ب$N^2$ أكبر منه في \omterm{def:b1:fluid-envelopes-structure:ocean-layers}{الترموكلاين} بنحو خمس مرّات.
  \item عند G، $\theta = 214.15 \times (1000/186)^{2/7} = \qty{346}{K}$: \omterm{def:b1:fluid-envelopes-structure:potential-temperature}{درجة الحرارة الكامنة} عند التروبوبوز نحو \qty{346}{K}.
\end{enumerate}
\end{solution}
